% GENERATED FILE. Edit canonical lessons, then run npm run generate:books.
% canonical-source-hash: fnv1a64:77c155d6660ed44d
% canonical-lessons: MR-C71-pot, MR-C71-khanda, MR-C71-bot, MR-C71-man, MR-C71-chehra

\chapter{Stomach, Shoulder, Finger, Neck, Face}
\label{ch:mr-stomach-shoulder-finger-neck-face}
\hlchaptermodality{71}

\begin{chapteropening}
\textbf{By the end of this chapter:} \emph{I can say the stomach, a shoulder, a finger, the neck and the face.}

Complete the last lesson of chapter 71: I can say the stomach, a shoulder, a finger, the neck and the face.
\end{chapteropening}

\section[\texorpdfstring{\mr{पोट} (poṭ)}{पोट (poṭ)}]{\mr{पोट} (poṭ) — the stomach}

\label{lesson:MR-C71-pot}



\begin{quote}
Before the new one: say the Marathi for a tooth, then the Marathi for the throat.
\end{quote}

\subsection*{You'll want to know: \mr{पोट}}
\textbf{\mr{पोट}} — \emph{poṭ} — \textquotedblleft{}the stomach\textquotedblright{}.

\begin{cousinweb}
\textbf{\mr{पोट}} (\emph{poṭ}) goes back to \emph{*pōṭṭa}, \textquotedblleft{}belly\textquotedblright{}, one of the forms CDIAL reconstructs under \emph{*pēṭṭa}, \textquotedblleft{}belly\textquotedblright{} (CDIAL 8376). CDIAL notes that these may be the same as its words for \textquotedblleft{}basket\textquotedblright{}.
\end{cousinweb}

\begin{grammarlens}[title={What you've built}]
One more word for this chapter.
\end{grammarlens}

\subsection*{Your turn}
\begin{itemize}
\raggedright
  \item \emph{Say it:} \emph{poṭ}
  \item \emph{Say it:} \emph{poṭ}, once more
  \item \emph{Say it:} say \emph{poṭ}
  \item \emph{Recall:} say the Marathi for a tooth, then the Marathi for the throat, then say \emph{poṭ} again
\end{itemize}

\begin{tcolorbox}[breakable,colback=teal!4,colframe=teal!35!black,title={Before you move on}]
What does \mr{पोट} mean? (\textquotedblleft{}The stomach\textquotedblright{}.) Say it once more.
\end{tcolorbox}
\section[\texorpdfstring{\mr{खांदा} (khāndā)}{खांदा (khāndā)}]{\mr{खांदा} (khāndā) — a shoulder}

\label{lesson:MR-C71-khanda}



\begin{quote}
Before the new one: say the Marathi for the throat, then the Marathi for the stomach.
\end{quote}

\subsection*{You'll want to know: \mr{खांदा}}
\textbf{\mr{खांदा}} — \emph{khāndā} — \textquotedblleft{}a shoulder\textquotedblright{}.

\begin{cousinweb}
\textbf{\mr{खांदा}} (\emph{khāndā}) is a tadbhava, a word inherited from Sanskrit and reshaped by sound change. It comes from \emph{skandha}, \textquotedblleft{}shoulder\textquotedblright{} (CDIAL 13627). CDIAL glosses the Marathi word \textquotedblleft{}shoulder, back of neck; large bough\textquotedblright{}.
\end{cousinweb}

\begin{grammarlens}[title={What you've built}]
One more word for this chapter.
\end{grammarlens}

\subsection*{Your turn}
\begin{itemize}
\raggedright
  \item \emph{Say it:} \emph{khāndā}
  \item \emph{Say it:} \emph{khāndā}, once more
  \item \emph{Say it:} say \emph{khāndā}
  \item \emph{Recall:} say the Marathi for the throat, then the Marathi for the stomach, then say \emph{khāndā} again
\end{itemize}

\begin{tcolorbox}[breakable,colback=teal!4,colframe=teal!35!black,title={Before you move on}]
What does \mr{खांदा} mean? (\textquotedblleft{}A shoulder\textquotedblright{}.) Say it once more.
\end{tcolorbox}
\section[\texorpdfstring{\mr{बोट} (boṭ)}{बोट (boṭ)}]{\mr{बोट} (boṭ) — a finger}

\label{lesson:MR-C71-bot}



\begin{quote}
Before the new one: say the Marathi for the stomach, then the Marathi for a shoulder.
\end{quote}

\subsection*{You'll want to know: \mr{बोट}}
\textbf{\mr{बोट}} — \emph{boṭ} — \textquotedblleft{}a finger\textquotedblright{}.

\begin{cousinweb}
The origin of \textbf{\mr{बोट}} (\emph{boṭ}) is not settled. Neither CDIAL nor DEDR, searched by form and by meaning, has an entry for it.
\end{cousinweb}

\begin{grammarlens}[title={What you've built}]
One more word for this chapter.
\end{grammarlens}

\subsection*{Your turn}
\begin{itemize}
\raggedright
  \item \emph{Say it:} \emph{boṭ}
  \item \emph{Say it:} \emph{boṭ}, once more
  \item \emph{Say it:} say \emph{boṭ}
  \item \emph{Recall:} say the Marathi for the stomach, then the Marathi for a shoulder, then say \emph{boṭ} again
\end{itemize}

\begin{tcolorbox}[breakable,colback=teal!4,colframe=teal!35!black,title={Before you move on}]
What does \mr{बोट} mean? (\textquotedblleft{}A finger\textquotedblright{}.) Say it once more.
\end{tcolorbox}
\section[\texorpdfstring{\mr{मान} (mān)}{मान (mān)}]{\mr{मान} (mān) — the neck}

\label{lesson:MR-C71-man}



\begin{quote}
Before the new one: say the Marathi for a shoulder, then the Marathi for a finger.
\end{quote}

\subsection*{You'll want to know: \mr{मान}}
\textbf{\mr{मान}} — \emph{mān} — \textquotedblleft{}the neck\textquotedblright{}.

\begin{cousinweb}
\textbf{\mr{मान}} (\emph{mān}) is a tadbhava, a word inherited from Sanskrit and reshaped by sound change. It comes from \emph{manyā}, \textquotedblleft{}nape of the neck\textquotedblright{} (CDIAL 9858). CDIAL glosses the Marathi word \textquotedblleft{}neck\textquotedblright{}.
\end{cousinweb}

\begin{grammarlens}[title={What you've built}]
One more word for this chapter.
\end{grammarlens}

\subsection*{Your turn}
\begin{itemize}
\raggedright
  \item \emph{Say it:} \emph{mān}
  \item \emph{Say it:} \emph{mān}, once more
  \item \emph{Say it:} say \emph{mān}
  \item \emph{Recall:} say the Marathi for a shoulder, then the Marathi for a finger, then say \emph{mān} again
\end{itemize}

\begin{tcolorbox}[breakable,colback=teal!4,colframe=teal!35!black,title={Before you move on}]
What does \mr{मान} mean? (\textquotedblleft{}The neck\textquotedblright{}.) Say it once more.
\end{tcolorbox}
\section[\texorpdfstring{\mr{चेहरा} (chehrā)}{चेहरा (chehrā)}]{\mr{चेहरा} (chehrā) — the face}

\label{lesson:MR-C71-chehra}



\begin{quote}
Before the new one: say the Marathi for a finger, then the Marathi for the neck.
\end{quote}

\subsection*{You'll want to know: \mr{चेहरा}}
\textbf{\mr{चेहरा}} — \emph{chehrā} — \textquotedblleft{}the face\textquotedblright{}.

\begin{cousinweb}
\textbf{\mr{चेहरा}} (\emph{chehrā}) is a loan, usually traced to Persian \emph{chihra}, \textquotedblleft{}face, countenance\textquotedblright{}.
\end{cousinweb}

\begin{grammarlens}[title={What you've built}]
One more word for this chapter.
\end{grammarlens}

\subsection*{Your turn}
\begin{itemize}
\raggedright
  \item \emph{Say it:} \emph{chehrā}
  \item \emph{Say it:} \emph{chehrā}, once more
  \item \emph{Say it:} say \emph{chehrā}
  \item \emph{Recall:} say the Marathi for a finger, then the Marathi for the neck, then say \emph{chehrā} again
\end{itemize}

\begin{tcolorbox}[breakable,colback=teal!4,colframe=teal!35!black,title={Before you move on}]
What does \mr{चेहरा} mean? (\textquotedblleft{}The face\textquotedblright{}.) Say it once more.
\end{tcolorbox}
